\section{地方登记、边郡军费与东汉的社会裂缝}
\label{sec:eastern-han-regions-households-frontier}

光武重建后，雒阳朝廷必须重新取得田土、人户、道路和地方官的可见性。建武十五年清查垦田、户口的度田行动，说明国家试图把战争后的土地与人口重新放入可征收、可调发的框架；它也立即暴露了登记触及地方占有和豪强关系时会发生的摩擦。度田并不是一张完整普查表，诏令与史书不能证明每个郡县的丈量、隐漏和负担相同。它更适合被看作东汉国家以文书面对地区差异的一次努力。

\subsection{交趾与北边：行政名目之外的协商}

建武十六年徵侧、徵贰起事，建武十九年马援平定交趾。起事、征讨和大体结果可以由《后汉书》帝纪、马援传及交趾相关材料相互核对；地方参与者的目的、性别、族属和战争损失却不能由胜者叙事一概裁定。岭南的河流、海路、县治、地方首领和税役网络，同雒阳所见并不相同。平定的意义在于朝廷重新安排地方军政与名义，不在于证明交趾社会从此被无缝纳入中原县乡。

建武二十四年南匈奴日逐王比归附，汉廷在边郡安置其部众。南匈奴既可成为北边缓冲和军事伙伴，也保有自身首领和资源网络；“归附”并不等于部众数量、驻牧界线或自主程度已可由朝廷单方面决定。七三年东汉在伊吾卢取得据点、班超赴鄯善，七五年陈睦在车师危机中遇害、朝廷一度撤出，这一推进与后撤更能说明西域经营依赖河西补给、绿洲合作与道路安全。都护的名义不是直接统治的同义词，驻军和使节关系必须不断被物资、外交和地方王权维持。

\subsection{羌乱如何变成财政与迁徙问题}

七六年陇西、金城等地的羌乱加剧，一〇七年凉州、并州再起，一六七年至一六九年段颎高强度进军，构成延续数十年的边郡压力。史书可保存朝廷出师、将领任命和胜败序列，却极少留下地方粮秣如何筹集、被迁徙者如何安置、军户如何维持的连续账簿。段颎的“平定”短期恢复了部分通道与郡县控制，却未消除征发、牧地、地方联盟和军费的条件；一八四年凉州军民与羌胡部众再起，正说明暴力的结局不等于地方秩序的结局。

一七七年鲜卑诸部在檀石槐领导下侵扰北边，汉军出击受挫。北方威胁的变化同西部羌乱叠加，使募兵、军费和地方协力成为雒阳难以回避的负担。“鲜卑”“羌胡”都是当时官府、史家在不同情境中使用的分类，不可被写成内部统一、目标一致的社会实体。把边患还原为多种部众、边郡吏员、军人和地方居民的关系，才不会将东汉后期的财政困局误读为单一民族冲突。

\subsection{学习、宗教与地方声望的并行网络}

永平八年楚王英案中的黄老与佛教供养用语，七九年白虎观讨论经义，一七五年熹平石经立于太学，一七八年鸿都门学招纳辞赋书画之士，显示国家不断试图规定何种知识、仪式和文字能够取得公开权威。经学会议、石经和学校并不代表地方教育已经同质；鸿都门学也不只是“非儒”的反面。不同文本、书写技巧和师友关系都可成为接近宫廷、人事和地方声望的途径。石经残石、传本和太学遗址能校验部分物质事实，却不能把学校制度外推为所有读书人都经历的生活。

太平道在一八三年前后组织信众，黄巾在一八四年扩展，说明宗教语言、互助、地方流动和政治动员可以相互转换。材料多来自平叛者和后来史家，人数、组织层级和各地联络程度均不可照录；黄巾主力在广宗、下曲阳受挫，也不等于村落、佃作和逃户问题已经消失。宗教与知识网络之所以重要，不是它们天然反国家，而是它们提供了县廷、豪强、师友和军队之外的联结方式。

\subsection{州牧、屯田与皇帝名义的再分配}

一八八年置州牧、扩大州刺史军政权限，是在黄巾、凉州与边郡多重危机中的应对。制度存在可靠，覆盖地域、权力大小和后来的割据关系却因州而异；不能将三国的结果倒投为当时唯一目的。幽州张纯、张举的起事与乌桓关系的动荡，又显示东北地方军人可以利用中央分身乏术的时机组织兵力。州牧、郡守、豪强和边部之间谁能掌握粮仓、兵员和道路，已比一份任命书更直接地影响实际控制。

一九六年曹操迎献帝至许县并组织屯田，将皇帝诏令、兖州粮仓、军队和行政人员编入同一军政框架；二〇七年北征乌桓，控制河北与幽州通道。屯田初创时间、规模和个别人物的作用有争议，但它清楚表明：皇帝名义只有同供给和文书结合才会产生持续能力。东汉的终局不是某一日被废除的制度，而是度田、边防、学校、宗教、州牧和屯田这些原本可以联结地方与中枢的机制，陆续被不同军事集团占有和重组。

\subsection{史料的缺口}

《后汉书》帝纪、度田、马援、班超、西羌、鲜卑、灵帝与献帝相关传记，《后汉纪》及《三国志》共同构成编年主干；石经、边塞遗址、地方志性材料和后出的宗教文本提供不同尺度。军功、斩获、部众称谓、户口数和“归附”程度尤其需要保留材料立场。东汉社会并不是正史中某一次政变的附属背景，而是由登记、迁徙、军费、知识与地方中介持续塑形的多区域世界。
